% notes-definability.tex
% Topic: Definability

\begin{definitionbox}{Definition (Definable Set)}
\begin{definition}[Definable Set]\label{def:definable-set}
Let $\mathcal{M}$ be an $\mathcal{L}$-structure. A set
$X\subseteq M^n$ is \textbf{definable in $\mathcal{M}$} if there is an
$\mathcal{L}$-formula $\varphi(x_1,\ldots,x_n)$, possibly with parameters
from $M$, such that
\[
  X=\{\bar a\in M^n : \mathcal{M}\models\varphi(\bar a)\}.
\]
\end{definition}
\end{definitionbox}

\begin{remark*}[Standard quantified statement]
$X$ is definable in $\mathcal{M}$ iff
$\exists\varphi(x_1,\ldots,x_n)$ such that
$\bar a\in X$ iff $\mathcal{M}\models\varphi(\bar a)$.
\end{remark*}

\begin{remark*}[Predicate reading]
\[
\mathsf{DefinableSet}(X,\mathcal{M})
\Longleftrightarrow
\exists\varphi\ \forall\bar a\,
\bigl(\bar a\in X \leftrightarrow \mathcal{M}\models\varphi(\bar a)\bigr).
\]
\end{remark*}

\begin{remark*}[Interpretation]
Definable sets are the subsets of powers of the universe that can be described
inside the language of the structure.
\end{remark*}

\begin{dependencies}
\begin{itemize}
  \item \hyperref[def:l-structure]{$\mathcal{L}$-Structure}
  \item \hyperref[def:l-formula]{$\mathcal{L}$-Formula}
  \item \hyperref[def:formula-interpretation]{Interpretation of Formulas}
\end{itemize}
\end{dependencies}

\begin{definitionbox}{Definition (Definable Function)}
\begin{definition}[Definable Function]\label{def:definable-function}
Let $\mathcal{M}$ be an $\mathcal{L}$-structure. A function
$g\colon X\to M^k$, with $X\subseteq M^n$, is \textbf{definable in
$\mathcal{M}$} if its graph
\[
  \Gamma_g=\{(\bar a,\bar b)\in M^{n+k}: \bar a\in X \text{ and } g(\bar a)=\bar b\}
\]
is a definable set in $\mathcal{M}$.
\end{definition}
\end{definitionbox}

\begin{remark*}[Standard quantified statement]
$g$ is definable in $\mathcal{M}$ iff its graph $\Gamma_g$ is definable in
$\mathcal{M}$.
\end{remark*}

\begin{remark*}[Interpretation]
A definable function is one whose input-output relation can be described by a
formula in the structure.
\end{remark*}

\begin{dependencies}
\begin{itemize}
  \item \hyperref[def:definable-set]{Definable Set}
\end{itemize}
\end{dependencies}

\subsection{Basic Examples and Closure}

\begin{remark*}[Domain and image]
If the graph of a function is defined by a formula
$\varphi(\bar x,\bar y)$, then its domain is defined by
$\exists\bar y\,\varphi(\bar x,\bar y)$ and its image is defined by
$\exists\bar x\,\varphi(\bar x,\bar y)$.
\end{remark*}

\begin{example*}[Absolute value in an ordered ring]
In the ordered-ring language, the absolute-value function on the real field is
definable by the formula
\[
  \bigl((x<0)\land (y=-x)\bigr)
  \lor
  \bigl(\lnot(x<0)\land (y=x)\bigr).
\]
This formula defines the graph of the function $x\mapsto |x|$.
\end{example*}

\begin{lemma}[Closure of Definable Sets]\label{lem:closure-of-definable-sets}
Let $\mathcal{A}$ be an $\mathcal{L}$-structure. If
$S_1,S_2\subseteq A^n$ are definable, then
$S_1\cap S_2$, $S_1\cup S_2$, and $A^n\setminus S_1$ are definable.
Furthermore, if $S\subseteq A^{n+m}$ is definable, then its projection
\[
  \mathrm{proj}(S)
  =
  \{\bar a\in A^n : \text{there is some }\bar b\in A^m
      \text{ such that }(\bar a,\bar b)\in S\}
\]
is definable.
\hyperref[prf:closure-of-definable-sets]{\textit{Go to proof.}}
\end{lemma}

\begin{remark*}[Standard quantified statement]
Definable subsets of $A^n$ are closed under finite Boolean operations, and
definable subsets of $A^{n+m}$ have definable coordinate projections.
\end{remark*}

\begin{remark*}[Predicate reading]
\[
\mathsf{DefinableSet}(S,\mathcal{A})
\Longrightarrow
\mathsf{DefinableSet}(\mathrm{proj}(S),\mathcal{A}).
\]
\end{remark*}

\begin{remark*}[Interpretation]
Boolean connectives give the Boolean closure operations on definable sets, and
existential quantification corresponds to projection.
\end{remark*}

\begin{dependencies}
\begin{itemize}
  \item \hyperref[def:definable-set]{Definable Set}
  \item \hyperref[def:l-formula]{$\mathcal{L}$-Formula}
  \item \hyperref[def:formula-interpretation]{Interpretation of Formulas}
\end{itemize}
\end{dependencies}

\subsection{Projection Example}

\begin{example*}[Projection of a disc]
In an ordered-ring structure on the real field, the formula
\[
  (x-4)^2+(y-3)^2<4
\]
defines the open disc of radius $2$ centered at $(4,3)$ in the plane. The
formula
\[
  \exists y\,\bigl((x-4)^2+(y-3)^2<4\bigr)
\]
defines its projection to the $x$-axis, namely the interval $(2,6)$.
Subtraction and squaring may be treated as abbreviations when the relevant
functions are already definable in the structure.
\end{example*}

\subsection{Automorphism Preservation}

\begin{proposition}[Definable Sets Are Preserved Under Automorphisms]\label{prop:definable-sets-preserved-under-automorphisms}
Let $\mathcal{A}$ be an $\mathcal{L}$-structure. If
$S\subseteq A^n$ is definable and $\pi$ is an automorphism of $\mathcal{A}$,
then $\pi(S)=S$. In this case, $S$ is said to be preserved under all
automorphisms.
\hyperref[prf:definable-sets-preserved-under-automorphisms]{\textit{Go to proof.}}
\end{proposition}

\begin{remark*}[Standard quantified statement]
For every definable $S\subseteq A^n$ and every
$\pi\in\mathrm{Aut}(\mathcal{A})$, one has $\pi(S)=S$.
\end{remark*}

\begin{remark*}[Predicate reading]
\[
\mathsf{DefinableSet}(S,\mathcal{A})
\land
\mathsf{Automorphism}(\pi,\mathcal{A})
\Longrightarrow
\pi(S)=S.
\]
\end{remark*}

\begin{remark*}[Interpretation]
Automorphisms are invisible to formulas. Therefore a definable set cannot
distinguish between elements or tuples that are related by an automorphism.
\end{remark*}

\begin{dependencies}
\begin{itemize}
  \item \hyperref[def:definable-set]{Definable Set}
  \item \hyperref[def:l-structure-automorphism]{$\mathcal{L}$-Structure Automorphism}
  \item \hyperref[prop:isomorphisms-preserve-formulas]{Isomorphisms Preserve Formulas}
\end{itemize}
\end{dependencies}

\subsection{Nondefinability Example}

\begin{example*}[Order is not definable in the additive real group]
In the additive-group language, the map $\pi\colon\mathbb{R}\to\mathbb{R}$
given by $\pi(x)=-x$ preserves $0$ and addition:
\[
  \pi(x+y)=-(x+y)=(-x)+(-y)=\pi(x)+\pi(y).
\]
It is its own inverse, hence an automorphism of the additive real group.
However, $0<1$ while $\pi(0)\not<\pi(1)$. Thus the order relation is not
preserved by this automorphism, so it is not definable in the additive-group
structure on $\mathbb{R}$.
\end{example*}

\begin{lemma}[No Nontrivial Automorphisms of the Ordered Real Ring]\label{lem:no-nontrivial-automorphisms-ordered-real-ring}
There are no nontrivial automorphisms of the ordered-ring structure on
$\mathbb{R}$.
\hyperref[prf:no-nontrivial-automorphisms-ordered-real-ring]{\textit{Go to proof.}}
\end{lemma}

\begin{remark*}[Standard quantified statement]
If $\pi$ is an automorphism of the ordered-ring structure on $\mathbb{R}$,
then $\pi$ is the identity map.
\end{remark*}

\begin{remark*}[Predicate reading]
\[
\mathsf{Automorphism}(\pi,\mathbb{R}_{\mathrm{oring}})
\Longrightarrow
\pi=\mathrm{id}_{\mathbb{R}}.
\]
\end{remark*}

\begin{remark*}[Interpretation]
The automorphism test is useful for proving nondefinability only when the
structure has nontrivial automorphisms. In the ordered real field, the identity
is the only automorphism, so every subset is preserved by all automorphisms
even though most subsets are not definable.
\end{remark*}

\begin{dependencies}
\begin{itemize}
  \item \hyperref[def:l-structure-automorphism]{$\mathcal{L}$-Structure Automorphism}
  \item \hyperref[prop:definable-sets-preserved-under-automorphisms]{Definable Sets Are Preserved Under Automorphisms}
\end{itemize}
\end{dependencies}
